\section{A consequence for almost-embeddable families}
\label{sec:consequences}

Lee and Sidiropoulos reduce their fixed-parameter almost-embeddable
families to planar and bounded-treewidth pieces joined at single
vertices~\cite[Theorem~1.7]{LS09}. We combine this reduction with
Theorem~\ref{thm:main} and the bounded-treewidth embedding theorem
of the companion~\cite[Theorem~1.1]{Treewidth108}.

A $1$-sum joins two graphs by identifying one vertex of each.
A stochastic embedding of distortion $D$ is a distribution of maps
into graph metrics that do not contract distances, with expected
expansion at most $D$ for every pair.

Informally, the construction of these families starts with a graph drawn
on a compact surface of genus $g$, orientable or nonorientable. One adds
$k$ attachments of prescribed width $w$, called fringes, to $k$ faces,
and then $a$ additional vertices with arbitrary adjacencies; these last
vertices are the apices. In particular, when $k=0$ this includes
bounded-genus graphs with a bounded number of arbitrarily attached
vertices. The corollary uses Lee and Sidiropoulos's named family and
their prescribed attachment and width conventions.

\begin{corollary}[Fixed almost-embeddable families]\label{cor:almost-embeddable}
Fix nonnegative integers $g,k,w,a$, and interpret
$\mathrm{AE}(g,(k,w),a)$ exactly as the almost-embeddable family of
Lee and Sidiropoulos~\cite[Theorem~1.7]{LS09}. Let
$h_0=h_0(g,k,w,a)\in\N_0$ be the target treewidth bound and let
$D=D(g,k,w,a)\ge1$ be a finite stochastic distortion bound supplied there.
Put
\[
 \widehat h=\max\{h_0,1\},\qquad
 K(g,k,w,a)=D\max\{C_{\mathrm{pl}},C_{\widehat h+1}\},
\]
where $C_{\mathrm{pl}}$ is the constant of Theorem~\ref{thm:main}, and
$C_t$ is the companion's constant for bags of cardinality at most
$t\ge2$. Every finite connected simple graph $G$ in this named family,
with strictly positive real edge lengths, admits a map $f$ into a real
$L_1$ space satisfying
\[
 d_G(u,v)\le\|f(u)-f(v)\|_1
 \le K(g,k,w,a)d_G(u,v)\qquad(u,v\in V(G)).
\]
For the same underlying graph, $K(g,k,w,a)$ also bounds the undirected
fractional concurrent flow--cut gap $\phi_*/\lambda_*$ for the finite
nonnegative capacities and demands defined in Section~\ref{intro:flow},
whenever $\lambda_*>0$.
\end{corollary}

\begin{proof}
The singleton case is immediate. Otherwise, replace each input edge length
by the distance between its endpoints; this preserves the metric and gives
reduced positive lengths. Lee--Sidiropoulos Theorem~1.7 stochastically embeds
this metric with distortion at most $D$ into finite $1$-sums of planar graphs
and graphs of treewidth at most $h_0$. Since the companion indexes $C_t$ by
bag cardinality, enlarging $h_0$ to $\widehat h$ gives the applicable constant
$C_{\widehat h+1}$. Put $M=\max\{C_{\mathrm{pl}},C_{\widehat h+1}\}$.

For a connected target piece with nonnegative edge lengths, add $\varepsilon>0$
to every edge and apply Theorem~\ref{thm:main} or the companion theorem.
As $\varepsilon\downarrow0$, the shortest-path matrices converge, so the
corresponding $L_1$ distance matrices have a subsequential limit in the closed
finite cut cone, retaining exactly the bound $M$ and identifying zero-distance
pairs. The finite expected expansion
ensures that all images of the finite input lie in a single finite-distance
target component almost surely. Lemma~1.5 of~\cite{LS09} retains the bound
$M$ on finite $1$-sums, and its Lemma~1.2 composes this bound with the
stochastic distortion $D$, giving $K=DM$. Applying the same cut-cone limit
to input lengths, followed by the fixed-graph equality
of~\cite[Section~3, Theorem~3.2]{GNRS04}, gives the flow--cut conclusion.
\end{proof}

This corollary concerns only the cited fixed-parameter almost-embeddable
families. It does not establish closure under general clique sums of order
greater than one, and therefore does not prove the full GNRS conjecture.
